\documentclass[12pt,a4paper]{article}
\usepackage[utf8]{inputenc}
\usepackage[T1]{fontenc}
\usepackage{lmodern}
\usepackage[french]{babel}
\usepackage[a4paper,margin=2.3cm]{geometry}
\usepackage{amsmath,amssymb}
\usepackage{enumitem}
\setlist[itemize]{label=\textbullet,nosep}
\usepackage[most]{tcolorbox}
\usepackage{fancyhdr}
\pagestyle{fancy}
\fancyhf{}
\fancyhead[L]{\small Exercice de maison}
\fancyhead[R]{\small Suites géométriques}
\fancyfoot[C]{\small\thepage}
\renewcommand{\headrulewidth}{0.4pt}
\setlength{\parindent}{0pt}
\setlength{\parskip}{5pt}
\newcommand{\question}[1]{\par\medskip{\large\bfseries #1}\par\smallskip}

\begin{document}

\begin{center}
{\LARGE\bfseries Corrigé}\\[4pt]
{\large Exercice de maison : suite géométrique}
\end{center}

On considère la suite $(V_n)$ définie par $V_0 = 5$ et, pour tout entier naturel $n$,
$V_{n+1} = 2V_n$.

\question{1. Calcul de $V_1$ et $V_2$}
On utilise la relation $V_{n+1} = 2V_n$.

Pour calculer $V_1$, on remplace $n$ par $0$ : \quad $V_1 = 2 \times V_0 = 2 \times 5 = 10$.

Pour calculer $V_2$, on remplace $n$ par $1$ : \quad $V_2 = 2 \times V_1 = 2 \times 10 = 20$.

Donc $\boxed{V_1 = 10 \ \text{ et } \ V_2 = 20}$.

\question{2. Montrons que $(V_n)$ est une suite géométrique}
On a $V_0 = 5$ et, pour tout entier naturel $n$ :
\[
V_{n+1} = 2 \times V_n.
\]
Chaque terme est le produit du terme qui le précède par la constante $2$, qui ne dépend pas de $n$.
D'après la définition, $(V_n)$ est donc une suite géométrique.

Donc $(V_n)$ est la suite géométrique de \textbf{raison} $\boxed{q = 2}$ et de
\textbf{premier terme} $\boxed{V_0 = 5}$.

\question{3. Expression de $V_n$ en fonction de $n$}
$(V_n)$ est une suite géométrique de premier terme $V_0 = 5$ et de raison $q = 2$ (non nulle).
Donc, pour tout entier naturel $n$ :
\[
V_n = V_0 \times q^{\,n} = 5 \times 2^{\,n}.
\]
\emph{Vérification :} $V_1 = 5 \times 2^1 = 10$ et $V_2 = 5 \times 2^2 = 20$.

Donc, pour tout entier naturel $n$ : \quad $\boxed{V_n = 5 \times 2^{\,n}}$.

\question{4. Sens de variation de $(V_n)$}
$(V_n)$ est une suite géométrique de premier terme $V_0 = 5$ et de raison $q = 2$.

Comme le premier terme et la raison sont tous deux supérieurs à zéro, la suite $(V_n)$ est une
suite à termes positifs.

Comme $q = 2$ et $2 > 1$, alors la suite $(V_n)$ est \textbf{croissante}.

\question{5. Calcul de $S = V_0 + V_1 + \dots + V_9$}
$S$ est la somme de termes consécutifs de la suite géométrique $(V_n)$, de raison $q = 2 \neq 1$.
\begin{itemize}
  \item Premier terme : $V_0 = 5$.
  \item Nombre de termes : $9 - 0 + 1 = 10$.
  \item Raison : $q = 2$.
\end{itemize}
On a donc :
\[
S = V_0 \times \frac{1 - q^{\,10}}{1 - q}
  = 5 \times \frac{1 - 2^{10}}{1 - 2}
  = 5 \times \frac{1 - 1024}{-1}
  = 5 \times \frac{-1023}{-1}
  = 5 \times 1023.
\]
Donc $\boxed{S = 5\,115}$.

\end{document}
